% Solutions of Book 4, chapter 19 (the transport layer).
\section*{Chapter \ref{ch:b2:transport-layer} --- The Transport Layer: UDP and TCP}

\begin{solution}{exo:b2:transport-layer:1}
The sender sends 0, 1, 2, 3; the acks of 0 and 1 let it send 4 and 5. Packet 2 is lost, so the receiver discards 3,
4 and 5 and acknowledges 1 each time. The timer of packet 2 expires: the sender goes back and sends 2, 3, 4, 5 again.
All arrive; the acks of 2, 3 and 4 let it send 6, 7 and 8. The ack of 5 is lost, but the ack of 6 is \omterm{def:b2:transport-layer:seq}{cumulative} and
covers 5, so 9 follows. In all: 0 1 2 3 4 5 2 3 4 5 6 7 8 9. The lost acknowledgement costs nothing.
\end{solution}

\begin{solution}{exo:b2:transport-layer:2}
It delivers packets 0, 1, 2 and 3 correctly, then the stale copy of packet 0, which it had stored as packet 4. The
fifth item delivered is wrong, and the real packet 4, when it comes, will be taken as a duplicate. With $k = 2$, a
\omterm{def:b2:transport-layer:window}{window} of $2 = 2^{k-1}$ prevents it (\cref{thm:b2:transport-layer:sr-window}): the receiver's \omterm{def:b2:transport-layer:window}{window} after packets 0
and 1 is $\{2, 3\}$, and a repeated 0 falls in the \omterm{def:b2:transport-layer:window}{window} before, which is only acknowledged again.
\end{solution}

\begin{solution}{exo:b2:transport-layer:3}
\omcode{code/bachelor-2/19-transport-layer/listings/stop_and_wait.py}{27}{38}{The receiver of \omterm{def:b2:transport-layer:rdt}{stop-and-wait}}
A duplicate arrives because the sender did not get the acknowledgement of the original, lost or late. If the receiver
did not acknowledge it again, the sender would send it for ever.
\end{solution}

\begin{solution}{exo:b2:transport-layer:4}
$L/R = 12\,000 / 10^9$~s $= 12$~µs, so $U = 0.012 / (30 + 0.012) \approx 0.0004$: 0.04~\% of the link. The \omterm{def:b2:transport-layer:window}{window} that
fills it is $(\mathrm{RTT} + L/R)/(L/R) \approx 2\,501$ packets, about 3.75 MB in flight: the product of the rate and
the round trip, which a long fast link needs in its \omterm{def:b2:transport-layer:window}{windows}.
\end{solution}

\begin{solution}{exo:b2:transport-layer:5}
First sample: SRTT 800, RTTVAR 400, timeout $800 + 4 \times 400 = 2400$. Second: RTTVAR $= \frac34 400 + \frac14 0 =
300$, SRTT 800, timeout 2000. Third: RTTVAR $= \frac34 300 + \frac14 |800 - 1600| = 425$, SRTT $= \frac78 800 +
\frac18 1600 = 900$, timeout $900 + 1700 = 2600$. Fourth: RTTVAR $= \frac34 425 + \frac14 100 = 343.75$, SRTT $=
887.5$, timeout 2262.5. These are the values \code{rto.py} prints for this series, none clamped since all exceed 1000:
\omoutput[1-1,10-14]{code/bachelor-2/19-transport-layer/outputs/rto.txt}{The second series of samples.}
The third timeout exceeds
the 1600 just measured because the variation term grows at once with the surprise, so that the next sample of 1600
would not cause a needless retransmission.
\end{solution}

\begin{solution}{exo:b2:transport-layer:6}
When packet 7 is lost and packet 8 arrives, the receiver acknowledges 8. The sender reads acknowledgements as
\omterm{def:b2:transport-layer:seq}{cumulative}, so it believes 7 arrived, slides past it and never sends it again. The receiver, still waiting for
sequence number 7 (modulo 16), later accepts packet 23, which has that number, as if it were packet 7: the stream is
corrupted, and from then on its expectations and the sender's no longer meet. Only 520 items are delivered, wrong
ones among them. A \omterm{def:b2:transport-layer:window}{go-back-N} receiver must acknowledge the last packet it accepted in order.
\end{solution}

\begin{solution}{exo:b2:transport-layer:7}
\omcode{code/bachelor-2/19-transport-layer/listings/selective_repeat.py}{5}{33}{The selective-repeat sender}
Each packet sent schedules its own expiry; an expiry does nothing if the packet has been acknowledged or slid past
in the meantime. An acknowledgement marks one packet, and the base slides past every marked packet at its front. The
tests run it with $W = 2^{k-1}$ under loss and duplication, and with jitter when the sequence space is large.
\end{solution}

\begin{solution}{exo:b2:transport-layer:8}
Number packets by their absolute index; packet $i$ carries $i \bmod 2^k$. Let the receiver expect packet $e$.

\emph{Data.} The receiver accepted $e - 1$, sent at some time $t$; at $t$ the sender's \omterm{def:b2:transport-layer:window}{window} held $e - 1$, so its
base was at least $e - W$, and the base never decreases. The channel keeps the order, so every packet arriving now was
sent after $t$, and is at least $e - W$. It is at most $e + W - 1$, since the sender's base cannot pass $e$ (it moves
only on acknowledgements of accepted packets). A packet $i$ with $i \equiv e$ and $|i - e| \le W < 2^k$ is $e$ itself,
so the receiver accepts exactly the right packet.

\emph{Acknowledgements.} They arrive in the order sent, and each acknowledges the last packet accepted, so an arriving
acknowledgement $a$ satisfies $\text{base} - 1 \le a \le \text{next} - 1$, where next $-$ base $\le W$. Its offset $a -
\text{base}$ ranges over $W + 1$ values from $-1$ to $W - 1$, which are distinct modulo $2^k$ exactly when $W + 1 \le
2^k$; then the sender reads every acknowledgement correctly. So $W \le 2^k - 1$ suffices.

\emph{Necessity.} With $W = 2^k$, let packets $0$ to $2^k - 1$ all arrive and all their acknowledgements be lost. The
sender goes back and repeats packet 0; the receiver, expecting packet $2^k$, whose number is 0, accepts the old packet
as new.
\end{solution}

\begin{solution}{exo:b2:transport-layer:9}
\omcode{code/bachelor-2/19-transport-layer/listings/fast_retransmit.py}{6}{20}{Going back at the third duplicate
acknowledgement}
\omoutput{code/bachelor-2/19-transport-layer/outputs/fast_retransmit.txt}{5000 packets, 2~\% loss, timeout 300 ms.}
An acknowledgement for the packet before the base is a duplicate: packets after a loss keep arriving and keep being
answered with the same number. Three of them make the loss nearly certain, a round trip after it, instead of the
300 ms of the timer. The packets sent are the same; the transfer takes 24.7 s instead of 54.0.
\end{solution}

\begin{solution}{exo:b2:transport-layer:10}
In the steady state, a loss halves the \omterm{def:b2:transport-layer:window}{window} from $W$ to $W/2$, and it then grows by one per round trip until it
reaches $W$ again and the next loss comes: a tooth lasts $W/2$ round trips. The packets sent in a tooth are $\frac W2
+ (\frac W2 + 1) + \dots + W \approx \frac W2 \cdot \frac{3W}{4} = \frac38 W^2$, the number of round trips times the
average \omterm{def:b2:transport-layer:window}{window}. One packet in $1/p$ is lost, so $\frac38 W^2 = 1/p$ and $W = \sqrt{8/(3p)}$. The throughput is the
packets of a tooth over its duration:
\[
\frac{1/p}{(W/2)\,\mathrm{RTT}} = \frac{2}{pW\,\mathrm{RTT}} = \frac{1}{\mathrm{RTT}} \cdot \frac{2}{p}
\sqrt{\frac{3p}{8}} = \frac{1}{\mathrm{RTT}} \sqrt{\frac{3}{2p}}
\]
segments per second, and multiplying by the segment size MSS gives the law.
\end{solution}

\begin{solution}{exo:b2:transport-layer:11}
\omoutput{code/bachelor-2/19-transport-layer/outputs/exo_rtt_unfair.txt}{Round trips of 1 and 2 time units, shared
losses.}
A flow's rate is its \omterm{def:b2:transport-layer:window}{window} divided by its round trip. Flow 1's rate grows by 1 per time unit; flow 2's \omterm{def:b2:transport-layer:window}{window} grows
by 1 every 2 units, so its rate grows by $1/2$ per 2 units, $1/4$ per unit. Between two shared losses each rate climbs
from half its peak, so the peaks, and the averages, stand in the ratio of the growth rates, $4 : 1$: shares $0.8$ and
$0.2$, and $J = (0.8 + 0.2)^2 / (2 (0.64 + 0.04)) \approx 0.735$. Under AIMD with shared losses, a flow's share falls as
the square of its round trip.
\end{solution}

\begin{solution}{exo:b2:transport-layer:12}
Under MIMD every step multiplies $x_1$ and $x_2$ by the same factor, $c$ or $b$, so $x_1/x_2$ never changes: flows that
start unequal stay unequal in the same ratio, for ever. Under additive increase and additive decrease, every step adds
the same amount to both, so $x_1 - x_2$ never changes, and again the flows do not converge. AIMD works because one of
its two moves, the decrease, shrinks the difference and the other keeps it.
\end{solution}

\begin{solution}{pb:b2:transport-layer:1}
\textbf{1.} $8000$ bits at 8 Mbit/s: 1 ms. $U = 1/(100 + 1) \approx 1~\%$.

\textbf{2.} $10\,000 \times 101$ ms $\approx 1010$ s, about 17 minutes.

\textbf{3.} A cycle fails when the packet or its acknowledgement is lost, about 2~\% of the time, and then costs a
timeout of 300 ms instead of 101 ms. That adds about $0.02 \times 300 = 6$ ms per packet on average: some 6~\%, a
minute more.

\textbf{4.} When an acknowledgement is lost the sender repeats the packet, and the receiver must tell the copy from the
next packet. With one packet in flight and no reordering, the packet that arrives is either the one expected or the
one just accepted: one bit tells which.

\textbf{5.} A delayed old copy can arrive when the receiver expects the same bit again, two packets later, and be
accepted as new.

\textbf{6.} No acknowledgement can arrive before a round trip; a shorter timer would repeat every packet.

\textbf{7.} $(\mathrm{RTT} + L/R)/(L/R) = 101$ packets.

\textbf{8.} \omterm{def:b2:transport-layer:window}{Go-back-N}: $2^4 - 1 = 15$. \omterm{def:b2:transport-layer:window}{Selective repeat}: $2^3 = 8$. Both are far below 101, so the sequence numbers
limit the utilisation to about 15~\% and 8~\%.

\textbf{9.} \omterm{def:b2:transport-layer:window}{Go-back-N} resends the whole \omterm{def:b2:transport-layer:window}{window} behind a loss, up to 15 packets; \omterm{def:b2:transport-layer:window}{selective repeat} resends the lost
packet only.

\textbf{10.} For 5\,000 packets, \omterm{def:b2:transport-layer:window}{go-back-N} sent 13\,408 and took 76.3 s; \omterm{def:b2:transport-layer:window}{selective repeat} sent 5\,531 and took 36.7 s.

\textbf{11.} A fast link with a long round trip needs millions of bytes in flight, and an old segment can survive for
minutes; 32-bit numbers give 4 GiB of space, so that a number does not come round again while an old segment with it
may still exist, at ordinary speeds.

\textbf{12.} \omterm{def:b2:transport-layer:flow-control}{Flow control}: each acknowledgement advertises the free space in the receiver's buffer, and the sender
keeps no more than that in flight.

\textbf{13.} Doubling from 1, the \omterm{def:b2:transport-layer:window}{window} exceeds 101 after 7 round trips ($2^7 = 128$).

\textbf{14.} $\sqrt{3/0.02} \approx 12.2$ segments per round trip, $12.2 \times 8000$ bits per 0.1 s: about 1 Mbit/s.

\textbf{15.} By the losses: 1 Mbit/s against a link of 8.

\textbf{16.} $\sqrt{3/(2p)} \ge 101$ gives $p \le 3/(2 \times 101^2) \approx 1.5 \times 10^{-4}$, one loss in about
6800 packets.

\textbf{17.} A queue that overflows must be drained quickly, which only a large cut does; and the multiplicative
decrease is what makes the shares converge (\cref{thm:b2:transport-layer:aimd-fair},
\cref{exo:b2:transport-layer:12}).

\textbf{18.} 9.3~\% below: the \omterm{def:b2:transport-layer:window}{window} is only about 16 segments, so whole segments, and the round in which the loss
falls, weigh in the average; the law assumes a continuous \omterm{def:b2:transport-layer:window}{window}.

\textbf{19.} Equal shares, half the link each, on average.

\textbf{20.} The index is 0.947 at round 25 and 0.995 at round 50: between the two.

\textbf{21.} Both flows still add one per round and halve together at each loss: their sum oscillates, while their
difference has vanished. The index measures equality, not stability.

\textbf{22.} One fifth, by \cref{exo:b2:transport-layer:11}.

\textbf{23.} \omterm{def:b2:transport-layer:fairness}{Fairness} holds per connection: ten connections against one other get about ten shares in eleven.

\textbf{24.} The $\sqrt{3/(2p)}$ law, checked by simulation: within 3~\% for periodic losses once $p \le 0.001$
(2.5~\% and 1.7~\% at $p = 0.001$ and $0.0003$), and a few per cent above it with random losses.
\end{solution}
